\documentclass[11pt]{article}
\usepackage{amsmath,amssymb,amsthm}
\usepackage[T1]{fontenc}
\usepackage{amsmath,amssymb,amsthm}
\usepackage[margin=1in]{geometry}
\usepackage[hidelinks]{hyperref}

\newcommand{\ZZ}{\mathbb{Z}}
\newcommand{\HH}{\mathbb{H}}

\newtheorem{theorem}{Theorem}
\newtheorem{lemma}[theorem]{Lemma}
\newtheorem{proposition}[theorem]{Proposition}

\newcommand{\fptitle}{Review of Problem 5 -- Astra}
\newcommand{\reviewtag}{Reviewer 3}
\usepackage{fp-review}

\begin{document}
\maketitle

\section{Review}

In a ring or a field, the identity $x+y+z = 0$ implies the identity $x^3 + y^3 + z^3 - 3xyz = 0$. This submitted note provides a counterexample for its counterpart for hyperfields, in which the addition is possibly multi-valued. 

The counterexample presented in this note follows the ideas of Hall, Prenowitz, and Connes--Consani. Concretely, Prenowitz and Connes--Consani connected hyperfields with point-line arrangements, and Hall gave constructions of an infinite point-line arrangement from $\mathbb{Z}$. The proofs written down are all correct. However, the historical remarks are only at the end of Section 3 but not in the main body of the submission, and they are very brief, acknowledging only the work of Connes--Consani. The extension construction of the point-line arrangements (Section 4) is almost identical to the construction and proof in Hall, and the construction of a hyperfield from a projective plane (Section 5) is almost identical to the construction and proof in Connes--Consani. 

For publication in a journal, I recommend a major revision. In particular, the author should remove the proofs that are already written in the literature and only cite the corresponding papers instead. 


\paragraph{Correctness.}
\begin{itemize}
	\item  1: essentially flawless mathematical argument
	\item  0.5: minor errors that are easily corrected
	\item  0: blank solution, or acknowledged partial result
	\item  $-0.5$: major errors
	\item  $-1$: proposed proof makes no progress towards the result.
\end{itemize}
Note that points should be docked for wrong statements/ solutions.

\noindent Score: 1

\paragraph{Exposition.}
\begin{itemize}
	\item  1: well-written and clear solution.
	\item  0: argument that is possible to follow with serious effort
	\item  $-1$: very poor exposition
\end{itemize}

\noindent Score: -1

\paragraph{Attributions.}
\begin{itemize}
	\item  1: solutions contains fully appropriate attributions
		to related work and theorems being used.
	\item  0: moderate issues with missing attributions or hallucinated references
	\item  $-1$: blatant plagiarism or hallucinations
\end{itemize}

\noindent Score: -1

\newpage
\section{Submission with inline comments}


\begin{abstract}
The proposed implication is false for general hyperfields. We construct a
countably infinite counterexample whose multiplicative group is infinite
cyclic. The construction uses an explicitly defined infinite difference set.
All required incidence properties and all hyperfield axioms are verified
directly.
\end{abstract}

\section{Statement and conventions}

We will prove the following.

\begin{theorem}\label{thm:main}
There exists a hyperfield
\[
    \HH=\{0\}\sqcup\{t^n:n\in\ZZ\},
    \qquad t^m t^n=t^{m+n},
\]
in which every element is its own additive inverse and
\[
    0\in 1\boxplus t\boxplus t^6,
\]
but
\[
    0\notin
    1\boxplus t^3\boxplus t^{18}
      \boxplus t^7\boxplus t^7\boxplus t^7.
\]
Consequently, the implication in the question fails for
\[
    x=1,\qquad y=t,\qquad z=t^6.
\]
\end{theorem}

Here the symbols \(t^n\) are pairwise distinct formal symbols, and \(1=t^0\).
In particular, the integer \(0\), when used as an exponent or an incidence
point, must be distinguished from the hyperfield element \(0\).

For subsets \(A,B\) of a set equipped with hyperaddition, we use the convention
\[
    A\boxplus B
    =\bigcup_{\substack{a\in A\\ b\in B}}(a\boxplus b).
\]
Singleton sets are identified with their elements in such expressions.

We use the usual definition of a commutative hyperfield: its nonzero elements
form an abelian multiplicative group, zero is multiplicatively absorbing,
and hyperaddition is a commutative, associative, nonempty-valued operation
with neutral element \(0\), unique additive inverses, and reversibility
\[
    c\in a\boxplus b
    \quad\Longleftrightarrow\quad
    b\in c\boxplus(-a).
\]
Multiplication must distribute over hyperaddition as an equality of sets.

\review{The definition of hyperfields should be made clearer. }

The connection between projective incidence geometry and hyperfields
extending the Krasner hyperfield is developed in
\cite[Sections~3--4]{CC}. The construction below, including the verification
of the hyperfield axioms, is self-contained.

\section{An infinite difference set}

\begin{lemma}\label{lem:difference-set}
There exists an infinite set \(D\subseteq\ZZ_{\geq 0}\) such that
\begin{enumerate}
    \item
    \[
        D\cap\{0,1,\ldots,38\}=\{0,1,6,20,23,27\};
    \]
    \item every nonzero integer has a unique expression \(d-e\), where
    \(d,e\in D\).
\end{enumerate}
\end{lemma}

\begin{proof}
For a finite set \(A\subseteq\ZZ_{\geq 0}\), put
\[
    \Delta(A)=\{a-b:a,b\in A,\ a>b\}.
\]
Say that \(A\) has distinct positive differences if each element of
\(\Delta(A)\) arises from exactly one ordered pair \((a,b)\).

Begin with
\[
    A_0=\{0,1,6,20,23,27\}.
\]
Its positive differences are
\[
    1,3,4,5,6,7,14,17,19,20,21,22,23,26,27.
\]
These are fifteen distinct integers, accounting for all
\(\binom{6}{2}=15\) pairs. Thus \(A_0\) has distinct positive differences.

We first describe an extension step. Suppose that \(A\) is a finite set
containing \(A_0\), that \(A\) has distinct positive differences, and that
\(n\geq 1\) does not belong to \(\Delta(A)\). Set
\[
    M=\max A,\qquad u=2M+n+1,
    \qquad A'=A\cup\{u,u+n\}.
\]
All positive differences in \(A'\) occur in the following four sets:
\[
    \Delta(A),\qquad
    \{n\},\qquad
    \{u-a:a\in A\},\qquad
    \{u+n-a:a\in A\}.
\]
These sets are pairwise disjoint. Indeed, \(n\notin\Delta(A)\), while every
integer in either of the last two sets is at least
\[
    u-M=M+n+1>\max\{M,n\}.
\]
Every element of \(\Delta(A)\) is at most \(M\). Finally, an equality
\[
    u-a=u+n-b
\]
would imply \(b-a=n\), contrary to \(n\notin\Delta(A)\).
Within each of the last two sets, distinct elements of \(A\) give distinct
differences. Therefore \(A'\) again has distinct positive differences.
Moreover, \(n\in\Delta(A')\).

Now define \(A_n\) recursively for \(n\geq 1\). If
\(n\in\Delta(A_{n-1})\), let \(A_n=A_{n-1}\). Otherwise apply the extension
step with \(A=A_{n-1}\) and this value of \(n\). Set
\[
    D=\bigcup_{n\geq 0}A_n.
\]
Every positive integer \(n\) belongs to \(\Delta(A_n)\), hence is a
difference of two elements of \(D\). Such a representation is unique:
any two competing representations would have all their entries in some
finite \(A_N\), contradicting the distinct-difference property of \(A_N\).
Reversing the ordered pair proves existence and uniqueness for negative
integers as well.

At every extension step \(M\geq 27\), so
\[
    u=2M+n+1\geq 56>38.
\]
Thus no new element at most \(38\) is ever introduced, proving the first
assertion. Finally, \(D\) is infinite, since a finite set has only finitely
many differences.
\end{proof}

\begin{lemma}\label{lem:incidence}
For \(k\in\ZZ\), define
\[
    L_k=D+k=\{d+k:d\in D\}.
\]
Regard the integers as points and the sets \(L_k\) as lines. Then:
\begin{enumerate}
    \item every two distinct points lie on exactly one line;
    \item every two distinct lines meet in exactly one point;
    \item every line is infinite;
    \item integer translations preserve the collection of lines.
\end{enumerate}
\end{lemma}

\begin{proof}
Let \(p\neq q\) be integers. By Lemma~\ref{lem:difference-set}, there is a
unique ordered pair \(d,e\in D\) satisfying
\[
    p-q=d-e.
\]
Then
\[
    k=p-d=q-e
\]
is the unique integer for which \(p,q\in L_k\). This proves the first
assertion.

For \(k\neq \ell\), a point of \(L_k\cap L_\ell\) corresponds precisely to
an ordered pair \(d,e\in D\) satisfying
\[
    k+d=\ell+e,
    \qquad\text{equivalently}\qquad
    d-e=\ell-k.
\]
There is exactly one such ordered pair, so the intersection consists of
exactly one point. Since each \(L_k\) is infinite, this also shows that
different indices give different lines.

The third assertion follows from the infinitude of \(D\), and the fourth
from
\[
    L_k+s=L_{k+s}.
\]
\end{proof}

For distinct integers \(m,n\), write \(L(m,n)\) for their unique common
line.

\section{Construction and verification of the hyperfield}

Let
\[
    \HH=\{0\}\sqcup\{t^n:n\in\ZZ\},
\]
with multiplication
\[
    t^m t^n=t^{m+n},\qquad 0a=a0=0.
\]
For distinct nonzero elements \(a=t^m\) and \(b=t^n\), define their line by
\[
    \Lambda(a,b)=\{t^r:r\in L(m,n)\}\subseteq\HH^\times.
\]
The incidence properties of Lemma~\ref{lem:incidence} hold equally for these
lines in \(\HH^\times\).

Define hyperaddition by
\[
    0\boxplus a=a\boxplus 0=\{a\}
\]
and, for \(a,b\neq 0\), by
\begin{equation}\label{eq:addition}
    a\boxplus b=
    \begin{cases}
        \{0,a\},&a=b,\\[2mm]
        \Lambda(a,b)\setminus\{a,b\},&a\neq b.
    \end{cases}
\end{equation}

\begin{proposition}\label{prop:hyperfield}
These operations make \(\HH\) a commutative hyperfield. Its additive
negation is the identity map.
\end{proposition}

\begin{proof}
The nonzero elements form an infinite cyclic multiplicative group with
identity \(1=t^0\). Zero is absorbing by definition.

\medskip
\noindent
\textbf{Nonempty sums, commutativity, and the neutral element.}
Every line is infinite, so all sums in \eqref{eq:addition} are nonempty.
Commutativity and the additive neutral-element axiom follow directly from
the definition.

\medskip
\noindent
\textbf{Additive inverses and reversibility.}
For all \(a,b\in\HH\), the definition gives
\begin{equation}\label{eq:inverse}
    0\in a\boxplus b
    \quad\Longleftrightarrow\quad a=b.
\end{equation}
Consequently, every element has a unique additive inverse, namely
\[
    -a=a.
\]

To verify reversibility, first suppose that \(a,b,c\) are all nonzero.
Then \(c\in a\boxplus b\) holds precisely when either all three elements
coincide, or they are pairwise distinct and lie on a common line.
Both conditions are symmetric in \(a,b,c\).
If at least one element is zero, membership holds precisely when the
other two elements are equal, again a symmetric condition. Hence
\[
    c\in a\boxplus b
    \quad\Longleftrightarrow\quad
    b\in c\boxplus a
    \quad\Longleftrightarrow\quad
    b\in c\boxplus(-a),
\]
as required.

\medskip
\noindent
\textbf{Associativity.}
If one of the three elements is zero, associativity follows from the
neutral-element rule. We therefore consider nonzero elements.

If all three elements equal \(a\), then
\[
    (a\boxplus a)\boxplus a
    =\{0,a\}\boxplus a
    =\{0,a\}.
\]

Suppose next that the multiset of the three elements is
\(\{a,a,b\}\), where \(a\neq b\), and let
\(\Lambda=\Lambda(a,b)\). Then
\[
    (a\boxplus a)\boxplus b
    =\{b\}\cup\bigl(\Lambda\setminus\{a,b\}\bigr)
    =\Lambda\setminus\{a\}.
\]
Also, uniqueness of the line through two distinct points gives
\[
\begin{split}
    (a\boxplus b)\boxplus a
    &=
    \bigcup_{s\in\Lambda\setminus\{a,b\}}
        \bigl(\Lambda\setminus\{s,a\}\bigr)\\
    &=\Lambda\setminus\{a\}.
\end{split}
\]
For the last equality, given \(v\in\Lambda\setminus\{a\}\), choose
\(s\in\Lambda\setminus\{a,b,v\}\), which is possible because
\(\Lambda\) is infinite.

Now suppose that \(a,b,c\) are pairwise distinct and lie on one line
\(\Lambda\). We claim that
\begin{equation}\label{eq:collinear}
    (a\boxplus b)\boxplus c=\{0\}\cup\Lambda.
\end{equation}
Every summand on the left is contained in \(\{0\}\cup\Lambda\).
Choosing the intermediate element \(c\in a\boxplus b\) produces both
\(0\) and \(c\). Given \(v\in\Lambda\setminus\{c\}\), choose
\[
    s\in\Lambda\setminus\{a,b,c,v\}.
\]
Then \(s\in a\boxplus b\) and \(v\in s\boxplus c\), proving
\eqref{eq:collinear}.

Finally, suppose that \(a,b,c\) are noncollinear. In this case
\begin{equation}\label{eq:noncollinear}
\begin{split}
    (a\boxplus b)\boxplus c
    =\HH^\times\setminus
    \bigl(
       \Lambda(a,b)\cup\Lambda(a,c)\cup\Lambda(b,c)
    \bigr).
\end{split}
\end{equation}
To prove the forward inclusion, take
\[
    s\in\Lambda(a,b)\setminus\{a,b\}.
\]
Since \(c\notin\Lambda(a,b)\), the elements \(s,c\) are distinct.
The line \(\Lambda(s,c)\) meets \(\Lambda(a,b)\) only at \(s\), and
meets each of \(\Lambda(a,c)\) and \(\Lambda(b,c)\) only at \(c\).
Therefore
\[
    s\boxplus c=\Lambda(s,c)\setminus\{s,c\}
\]
avoids all three lines on the right of \eqref{eq:noncollinear}, as well
as the hyperfield zero.

Conversely, let \(v\) lie outside those three lines.
The line \(\Lambda(c,v)\) meets \(\Lambda(a,b)\) in a unique point \(s\).
This point is neither \(a\) nor \(b\), since otherwise \(v\) would belong
to \(\Lambda(a,c)\) or \(\Lambda(b,c)\). Also \(v\neq s,c\).
Thus
\[
    s\in a\boxplus b,\qquad v\in s\boxplus c,
\]
proving the reverse inclusion.

The formulas just established are independent of the order in which the
three elements are listed. Together with commutativity, they imply
\[
    (a\boxplus b)\boxplus c=a\boxplus(b\boxplus c)
\]
in every case. Associativity for subsets follows by taking unions over
their elements; hence finite hypersums can be parenthesized freely.

\medskip
\noindent
\textbf{Distributivity.}
Multiplication by \(t^k\) acts on nonzero elements by translating their
exponents by \(k\). By Lemma~\ref{lem:incidence}, it carries lines to
lines, and in particular
\[
    t^k\Lambda(t^m,t^n)
    =\Lambda(t^{m+k},t^{n+k})
    \qquad(m\neq n).
\]
It follows from \eqref{eq:addition}, including its equal-argument and
zero-argument cases, that
\[
    t^k(a\boxplus b)=t^ka\boxplus t^kb.
\]
For multiplication by zero, both sides are \(\{0\}\), since
\(a\boxplus b\) is nonempty. Thus distributivity holds for all scalars.
All hyperfield axioms have now been verified.
\end{proof}

\section{Failure of the cubic implication}

\begin{proof}[Proof of Theorem~\ref{thm:main}]
Use the hyperfield constructed above and set
\[
    x=1,\qquad y=t,\qquad z=t^6.
\]
The integers \(0,1,6\) all belong to \(D\). Consequently,
\[
    t^6\in\Lambda(1,t)\setminus\{1,t\}=1\boxplus t.
\]
It follows that
\[
    0\in t^6\boxplus t^6
      \subseteq (1\boxplus t)\boxplus t^6.
\]
Thus the hypothesis in the question is satisfied.

Consider the line
\[
    \Lambda=\{t^r:r\in D-20\}.
\]
Because \(20,23,27\in D\), this line contains
\[
    1,\qquad t^3,\qquad t^7.
\]
On the other hand,
\[
    t^{18}\notin\Lambda,
\]
because \(18\in D-20\) would imply \(38\in D\), contrary to
Lemma~\ref{lem:difference-set}.

Write
\[
    u=1,\qquad v=t^3,\qquad w=t^{18},\qquad q=t^7,
    \qquad A=u\boxplus v\boxplus w.
\]
The points \(u,v,w\) are noncollinear, and
\(\Lambda(u,v)=\Lambda\). Formula~\eqref{eq:noncollinear} therefore gives
\[
    A\subseteq\HH^\times\setminus\Lambda.
\]
In particular,
\begin{equation}\label{eq:excluded}
    0\notin A,\qquad q\notin A.
\end{equation}

The repeated sum of \(q\) satisfies
\[
    q\boxplus q\boxplus q
    =\{0,q\}\boxplus q
    =\{0,q\}.
\]
Hence
\[
\begin{split}
    u\boxplus v\boxplus w\boxplus q\boxplus q\boxplus q
    &=A\boxplus\{0,q\}\\
    &=A\cup\bigcup_{a\in A}(a\boxplus q).
\end{split}
\]
The first set on the last line does not contain zero by
\eqref{eq:excluded}. By \eqref{eq:inverse}, a set \(a\boxplus q\)
contains zero exactly when \(a=q\); this is also excluded by
\eqref{eq:excluded}. Therefore
\[
    0\notin
    1\boxplus t^3\boxplus t^{18}
      \boxplus t^7\boxplus t^7\boxplus t^7.
\]

Finally, multiplication and negation in this hyperfield give
\[
    x^3=1,\qquad y^3=t^3,\qquad z^3=t^{18},
    \qquad -xyz=xyz=t^7.
\]
Thus the last displayed nonmembership is precisely the failure of the
conclusion asked about in the question.
\end{proof}

\begin{thebibliography}{9}

\bibitem{CC}
A.~Connes and C.~Consani,
\emph{The hyperring of ad\`ele classes},
Journal of Number Theory \textbf{131} (2011), no.~2, 159--194.

\end{thebibliography}

\end{document}
